\chapter{Lexicon}

\section{Keywords}

\code{if}, \code{do}, \code{else}, \code{when}, \code{for}, \code{in},
\code{repeat}, \code{break}, \code{skip}, \code{return}\\*
\code{default}, \code{import}, \code{module}\\*
\code{const}, \code{var}, \code{alias},
\code{self}, \code{function}, \code{tuple}, \code{record}, \code{struct},
\code{type}, \code{and}, \code{or}, \code{not}, \code{as}

\paragraph{Reserved keywords} Some keywords are reserved for future-proofing
to ensure backward compatibility in the eventuality they get used.

\code{def}, \code{let}, \code{ref}, \code{at}, \code{with}, \code{interface},
\code{enum}

\section{Operators}

\begin{table}[H]
\resizebox{\linewidth}{!}{
	\begin{tabular}{rcl}
		\textbf{Operator}               & \textbf{Token}  & \textbf{Description}                            \\\hline
		PLUS                            & \texttt{+}      & Arithmetic sum, concatenation                   \\\hline
		MINUS                           & \texttt{-}      & Arithmetic subtraction, intersection            \\\hline
		MULTIPLICATION                  & \texttt{*}      & Arithmetic multiplication                       \\\hline
		DIVISION                        & \texttt{/}      & Arithmetic division                             \\\hline
		MODULO                          & \texttt{\%}     & Arithmetic remainder                            \\\hline
		EXPONENTIAL                     & \texttt{\^}     & Arithmetic power, text repetition               \\\hline
		EQUAL                           & \texttt{=}      & Loose equality                                  \\\hline
		NOT EQUAL                       & \texttt{!=}     & Loose inequality                                \\\hline
		LESS EQUAL                      & \texttt{<=}     & Arithmetic and lexicographical comparison       \\\hline
		LESS                            & \texttt{<}      & Arithmetic and lexicographical comparison       \\\hline
		GREATER                         & \texttt{>}      & Arithmetic and lexicographical comparison       \\\hline
		GREATER EQUAL                   & \texttt{>=}     & Arithmetic and lexicographical comparison       \\\hline
		CONTAINS                        & \texttt{=\~{}}  & Checks if a set of elements contains an element \\\hline
		IS CONTAINED                    & \texttt{\~{}=}  & Reverse of CONTAINS                             \\\hline
		\multirow{2}{*}{BOOL NOT}       & \texttt{not}    & \multirow{2}{*}{Boolean logic negation}         \\
	                                	& \texttt{!}      &                                                 \\\hline
		\multirow{2}{*}{BOOL AND}       & \texttt{and}    & \multirow{2}{*}{Boolean logic and}              \\
	                                	& \texttt{\&}     &                                                 \\\hline
		\multirow{2}{*}{BOOL OR}        & \texttt{or}     & \multirow{2}{*}{Boolean logic or}               \\
	                                	& \texttt{|}      &                                                 \\\hline
		BITWISE NOT                     & \texttt{!!}     & Bitwise negation                                \\\hline
		BITWISE AND                     & \texttt{\&\&}   & Bitwise and                                     \\\hline
		BITWISE OR                      & \texttt{||}     & Bitwise or                                      \\\hline
		BITWISE XOR                     & \texttt{|||}    & Bitwise exclusive or                            \\\hline
		BITWISE LEFT SHIFT              & \texttt{\{<\}}  & Bitwise left shift                              \\\hline
		BITWISE RIGHT SHIFT             & \texttt{\{>\}}  & Bitwise right shift                             \\\hline
		BITWISE RIGHT UNSIGNED SHIFT    & \texttt{\{>>\}} & Bitwise right unsigned shift                    \\\hline
		IS                              & \texttt{is}     & Runtime type check, pattern matching            \\\hline
		AS                              & \texttt{as}     & Explicit type coercion                          \\\hline
	    \\
		\multicolumn{3}{c}{Pseudo-Operators}                                                                \\\hline
		DOT ACCESS                      & \texttt{.}      & Members access                                  \\\hline
		DOTREF ACCESS                   & \texttt{::}     & Function member access                          \\\hline
	\end{tabular}
}
\end{table}

\section{Other tokens}

\code{..}, \code{:}, \code{,}, \code{;}, \code{\#}, \code{(}, \code{)}, \code{[}, \code{]},
\code{\{}, \code{\}}, \code{<}, \code{>}, \code{=>}

\section{Comments}

There are two kinds of comments: single line and multiline.

\subsection{Single line comment}

\begin{minted}{antlr}
	comment: '#' comment_content? ('\n' | EOF) ;
	comment_content: [^\n]+ ;
\end{minted}

Every character after the comment token \code{\#} until a newline \code{\\n} is
considered part of the comment.

\subsection{Multiline comment}

\begin{minted}{antlr}
	comment_block: '#{' .* ('}#' | EOF) ;
\end{minted}

Every character between and the open comment token \code{\#\{} and the close
comment token \code{\}\#} is considered part of the comment.\\*
A multiline comment can also be inlined in the code.

\begin{minted}{\giovilexer}
if true #{ this is a comment }# do something()
\end{minted}

\paragraph{Conventions} By convention, \code{\#} is immediately followed by
a space, as the first character is reserved for "special" comments, for example
documentation: \code{\#? This is a documentation comment}

\todo{Specificare RFC riguardo language server e altro?}

\clearpage

\section{Literals}

\subsection{Boolean}

\begin{minted}{antlr}
	literal_boolean: 'true' | 'false' ;
\end{minted}

A \code{Bool} literal is simply expressed by \code{true} or \code{false}

\subsection{Numbers}

\begin{minted}{antlr}
	literal_number
		: literal_integer
		| literal_natural
		| literal_real
		;
	literal_integer
		: ( numeral_decimal
		  | numeral_hex
		  | numeral_octal
		  | numeral_binary
		  ) 'i'?
		;
	literal_natural
		: ( numeral_decimal
		  | numeral_hex
		  | numeral_octal
		  | numeral_binary
		  ) 'n'?
		;
	literal_real
		: decimal_floating_point_literal
		| hexadecimal_floating_point_literal
		| (numeral_hex | numeral_octal | numeral_binary) 'r'
		;
\end{minted}

Giovi supports three different type of numbers: integers, natural and real.

\paragraph{Suffixes} By default, a numeral is parsed as an integer; you can
add a suffix to the number to explicitate the type: \code{n} for naturals,
\code{i} for integers, \code{r} for reals. \code{69n} \code{69i} \code{69r}

\paragraph{Separators} To improve readability of large numbers, you can use
underscores to group digits; separators cannot appear at the beginning or the
end of the number. For example: \code{100\_000\_000}

\clearpage

\paragraph{Bases} Numeral literals may be expressed in multiple bases.\\*
By default, decimal (base 10) is used; to indicate a different base use the
appropriate prefix:
\begin{itemize}
	\item Decimal (base 10): \code{69420}
	\item Hexadecimal (base 16): \code{0x10F2C}
	\item Octal (base 8): \code{0o207454}
	\item Binary (base 2): \code{0b0001\_0000\_1111\_0010\_1100}
\end{itemize}

\paragraph{Scientific notation} For real numbers, the scientific notation is
supported: \code{1e69}, \code{6.9e-10}

\paragraph{Binary exponent} For hexadecimal real numbers, a binary exponent can
be used: \code{0x1p69}

\todo{semplificare?}

\clearpage

\subsection{Text}
\label{sec:text_syntax}

\begin{minted}{antlr}
	literal_normal_text
		: '"' ([^\n] | escape_sequence | interpolation)* '"'
		;

	literal_normal_text_block
		: '"{' (. | escape_sequence | interpolation)* '}"'
		;

	literal_exact_text
		: "'" ([^\n] | escape_sequence)* "'"
		;

	literal_exact_text_block
		: "'{" (. | escape_sequence)* "}'"
		;
\end{minted}

There are four types of \code{Text} literals (a.k.a. strings).

\begin{itemize}
	\item \textbf{Normal text}: delimited by double quotes \code{"}
	\item \textbf{Exact text}: delimited by single quptes \code{'}
	\item \textbf{Normal text block}: delimited by \code{"\{} and \code{\}"}
	\item \textbf{Exact text block}: delimited by \code{'\{} and \code{\}'}
\end{itemize}

\paragraph{Differences} All types of text literals can contain escape
sequences.\\*
Normal and exact text must be terminated on the same line (they cannot contain
a newline).\\*
Normal and exact text blocks can span over multiple lines and do not require
the escape of the text delimiter.\\*
Normal text and normal text block allow interpolation.

\paragraph{Escape sequences} Some special characters may not appear in the
text, unless when escaped by prefixing it with \code{$\backslash$} (backslash).

\begin{minted}{antlr}
	escape_sequence: character_escape | unicode_escape ;
	character_escape: '\' [btnfrs"'\\${}] ;
	unicode_escape: '\' 'u' HexDigit HexDigit HexDigit HexDigit ;
\end{minted}

\paragraph{Interpolation} Normal text and normal text blocks support
interpolation, allowing the "injecttion" of a variable or expression in the
text.\\*
To inject a variable, simply prefix it with \code{\$} (dollar): \code{"My name
is \$name"}.\\*
In order for an expression that is not a simple identifier to be injected, it
needs to be wrapped in braces: \code{"My name is \$\{user.getName()\}"}

\subsubsection{Text blocks}

Normal and exact text blocks allow the text to span over multiple lines,
without the need to escape the text delimiter.

\begin{listing}
	\label{lst:textblock_example}
	\begin{minted}{\giovilexer}
if true do {
	if true do {
		html Text <- "{
			<html>
				<head></head>
				<body>
					<p class="foo">Hello world</p>
				<body>
			</html>
			}"
	}
}
	\end{minted}
	\caption{Example of a textblock appearing indented in a nested scope}
\end{listing}

The last newline character will be included. If you don't want it, put the
closing tag on the previous line

\begin{minted}{\giovilexer}
t Text <- "{
	hello world}"
\end{minted}

\subparagraph{Incidental whitespaces} In a text block, white
spaces\footnotemark[1] are preserved, unless they're considered
"incidental". Incidental white spaces are those needed for code indentation.

\footnotetext[1]{
	White spaces are all those character that are "invisible", such as spaces,
	tabulations, newlines, formfeeds etc.
}

\begin{minted}[showtabs,tab=>,escapeinside=||]{\giovilexer}
if true do {
	if true do {
		html Text <- "{
			<html>
				<head></head>
				<body>
					<p>Hello world</p>
				<body>
			</html>
			}"
#^^^^^^^^^^^
# the first three tabulations are considered "incidental"
# and won't be included in the text
	}
}
\end{minted}

\clearpage

The amount of whitespaces that will be considered "incidental" is defined by
the line that has less of them.

\begin{listing}
	\label{lst:textblock_incidental_whitespaces}
	\caption{A textblock with its actual content highlighted}
	\begin{minted}[autogobble=false,showtabs,tab=>,escapeinside=||]{text}
			_out <- "{
			|\colorbox{gray!80}{<html>}|
			|\colorbox{gray!80}{	<head></head>}|
			|\colorbox{gray!80}{	<body>}|
			|\colorbox{gray!80}{		<p>Hello world</p>}|
			|\colorbox{gray!80}{	<body>}|
			|\colorbox{gray!80}{</html>}|
			}"

			_out <- "{
		|\colorbox{gray!80}{	<html>}|
		|\colorbox{gray!80}{		<head></head>}|
		|\colorbox{gray!80}{		<body>}|
		|\colorbox{gray!80}{			<p>Hello world</p>}|
		|\colorbox{gray!80}{		<body>}|
		|\colorbox{gray!80}{</html>}|
		|\colorbox{gray!80}{	\}"}|
	\end{minted}
\end{listing}

\nb{
	Only the highlighted part will actually be included in the text literal.\\*
	Note how line 16 \code{</html>} dictates the amount of incidental
	whitespaces for the second textblock.
}

\subparagraph{Indentation leader} Only one type of white space will be
considered as the indentation leader; in case of mixed types, the first
incidental whitespace of the first line (excluding the line with the opening
text block tag) will be considered the indentation leader character. So it is
advised to be consistent with code formatting to avoid ambiguities and
undesired side effects.

\todo{riferimento a convenzioni e best practices}

\begin{minted}[autogobble=false,showtabs,tab=>,escapeinside=||]{text}
			_out <- "{
		|\colorbox{gray!80}{	<html>}|
		|\colorbox{gray!80}{		<head></head>}|
		|\colorbox{gray!80}{		<body>}|
		|\colorbox{gray!80}{····		<p>Hello world</p>}|
		|\colorbox{gray!80}{		<body>}|
		|\colorbox{gray!80}{	</html>}|
		|\colorbox{gray!80}{	\}"}|
\end{minted}

Note how line 5 accidentally used four spaces instead of a tab; since the tab
was chosen as indentation leader, the spaces are considered non-incidental,
therefore they should be included in the text, hence line 5 now dictates the
amount of incidental whitespaces.

\subsection{Identifiers}

Any other word that is not a reserved keyword is considered an identifier.

\todo{
	TODO: completare

	TODO: restringere i caratteri permessi ad un subset unicode?
}

